\documentclass[11pt,a4paper]{article}
\usepackage{amsmath,amssymb,amsfonts,mathtools}
\usepackage{bm}
\usepackage{physics}
\usepackage{graphicx}
\usepackage{booktabs}
\usepackage{hyperref}
\usepackage{geometry}
\usepackage{enumitem}
\usepackage{cite}
\usepackage{xcolor}
\geometry{margin=1in}

\title{\textbf{Artificial Intelligence in Physics:\
From Data Analysis to Scientific Discovery}}

\author{}
\date{\today}

\begin{document}

\maketitle

\begin{abstract}
Artificial intelligence (AI), and in particular modern machine learning (ML), is
rapidly becoming part of the methodology of physics. Its role extends far beyond
the traditional use of computers for numerical calculation. Machine-learning
methods are now used for experimental data analysis, parameter estimation,
inverse problems, numerical simulation, surrogate modelling, experimental
design, control, theory construction, and the search for new physical laws.
Physics-informed machine learning further attempts to combine empirical data
with known physical principles, while neural operators provide a new approach
to learning mappings between functions and accelerating the solution of
parameterized differential equations.

This paper reviews the role of AI in physics from a methodological rather than
purely technological perspective. We distinguish between AI as a computational
tool, AI as an approximation to a known physical theory, and AI as a possible
instrument of scientific discovery. Particular attention is paid to the relation
between data-driven and theory-driven approaches. We argue that the central
problem is not whether AI can replace theoretical physics, but how machine
learning can be integrated with the existing hierarchy of physical concepts:
symmetry, conservation laws, variational principles, differential equations,
effective theories, and statistical inference.

The paper also discusses fundamental limitations. A machine-learning model can
fit data without identifying the underlying causal structure, can violate exact
physical constraints, can extrapolate catastrophically outside its training
domain, and can produce apparently convincing but scientifically meaningless
patterns. Consequently, validation, interpretability, uncertainty quantification,
dimensional analysis, symmetry, and independent physical verification remain
essential. The most promising future appears to be a hybrid methodology in
which human physicists and AI systems cooperate: physics supplies structure,
constraints, and questions, while AI supplies high-dimensional search,
approximation, pattern recognition, and computational exploration.
\end{abstract}

\tableofcontents

\section{Introduction}

Physics has always been intimately connected with computation. The development
of numerical methods transformed problems that could not be solved analytically
into quantitative predictions. Computational physics subsequently became a
third component of scientific methodology alongside experiment and theory.

Artificial intelligence introduces a further change. Traditional numerical
computation is based on algorithms explicitly constructed from mathematical
equations. Machine learning, in contrast, constructs an approximation from
examples, observations, simulations, or other forms of data. The distinction
can be expressed schematically as
\begin{equation}
\text{equations} \longrightarrow \text{algorithm}
\end{equation}
for conventional computational physics, whereas machine learning often follows
\begin{equation}
\text{data} \longrightarrow \text{model}.
\end{equation}

The distinction should not be exaggerated. Modern scientific machine learning
frequently combines the two:
\begin{equation}
\text{data}+\text{physical laws}+\text{numerical methods}
\longrightarrow
\text{learned physical model}.
\end{equation}
This hybrid approach is one of the most important developments in contemporary
computational science.

Machine learning has already become pervasive in physics, including
experimental particle physics, condensed matter physics, plasma physics,
fluid mechanics, cosmology, quantum physics, and climate physics
\cite{Carleo2019,Karagiorgi2022,Karn2021,Brunton2020}.
The application spectrum ranges from comparatively modest tasks, such as
classifying experimental events, to ambitious attempts to discover new
equations and physical structures.

The important question is therefore no longer whether AI can be useful in
physics. It clearly can. The more fundamental question is:

\begin{quote}
\emph{What is the proper epistemological and methodological role of AI in
physics?}
\end{quote}

This question is more subtle than the question of computational efficiency.
Physics is not merely the prediction of numbers. A physical theory attempts to
identify a coherent structure that explains a class of phenomena and permits
generalization to situations not previously observed. Prediction and
understanding are related, but they are not identical.

\section{AI, Machine Learning, and Physics}

The term artificial intelligence covers a large collection of methods. For
physics, three broad categories are especially important.

\subsection{Supervised learning}

In supervised learning one is given pairs $(x_i,y_i)\quad i=1\dots N$,
and attempts to construct a function $y\simeq f(\theta,x)$,where $\theta$ 
denotes the parameters of the model.

In physics, $x$ may represent experimental observables and $y$ may represent
particle identities, phases of matter, physical parameters, or solutions of
differential equations.

\subsection{Unsupervised and self-supervised learning}

Here the objective is to identify structure without specifying the desired
output explicitly. Examples include clustering, dimensional reduction,
anomaly detection, and representation learning.

This is particularly interesting for physics because scientific discovery often
begins with the recognition that data contain a structure not anticipated by
the existing theoretical description.

\subsection{Generative models}

Generative models attempt to represent a probability distribution
$p(x)$ or a conditional distribution $p(y|x)$.

Generative models are useful for simulation, sampling, uncertainty estimation,
and inverse problems. They can potentially replace expensive Monte Carlo or
numerical calculations by learned approximations.

Large language models represent another form of generative AI. Their immediate
importance for physics is not that they ``understand'' physical reality in the
same way that a physical theory does, but that they can manipulate large
amounts of scientific text, mathematical expressions, computer code, and
symbolic descriptions. Their role in theoretical research is therefore
potentially significant, but also requires special caution
\cite{He2024}.

\section{The Hierarchy of AI Applications in Physics}

It is useful to distinguish several levels of AI usage.

\begin{enumerate}[label=\arabic*.]
\item \textbf{Data analysis:} extracting information from experimental data.

\item \textbf{Parameter inference:} estimating physical parameters from
observations.

\item \textbf{Numerical acceleration:} replacing expensive calculations by
surrogate models.

\item \textbf{Inverse problems:} reconstructing hidden physical quantities from
observable consequences.

\item \textbf{Experimental design:} selecting measurements that maximize
scientific information.

\item \textbf{Model discovery:} finding mathematical relations compatible with
observations.

\item \textbf{Theory discovery:} constructing candidate physical principles or
dynamical laws.

\item \textbf{Scientific autonomy:} combining hypothesis generation, simulation,
experiment, analysis, and verification into an automated scientific loop.
\end{enumerate}

The first five levels are already established research activities. The last
three are substantially more ambitious and raise fundamental questions about
what constitutes scientific understanding.

\section{AI as a Tool for Experimental Physics}

\subsection{High-dimensional data}

Modern experiments generate enormous quantities of data. Particle detectors,
astronomical surveys, gravitational-wave observatories, synchrotron sources,
neutron facilities, and imaging systems all produce datasets whose dimensionality
can exceed the capacity of conventional analysis methods.

Machine learning is particularly effective when the useful information is
distributed nonlinearly over a large number of observables.

For example, a classifier can approximate
\begin{equation}
P(\mathrm{signal}|x_1,\ldots,x_n),
\end{equation}
where the $x_i$ are measured quantities.

In high-energy physics this can be used to distinguish signal from background.
Machine learning has consequently become an important component of searches for
new physics \cite{Karagiorgi2022}.

\subsection{Anomaly detection}

Anomaly detection is potentially more interesting from the point of view of
fundamental physics.

A conventional search begins with a hypothesis:
\begin{equation}
H_1=\text{specific new-physics model}.
\end{equation}

One then searches for the predicted signature.

An alternative is
\begin{equation}
H_0=\text{known physics}
\end{equation}
and asks whether the data contain statistically significant deviations from
the learned distribution of known events.

This changes the logic from
\begin{equation}
\text{theory}\rightarrow\text{prediction}\rightarrow\text{test}
\end{equation}
toward
\begin{equation}
\text{data}\rightarrow\text{anomaly}\rightarrow\text{hypothesis}.
\end{equation}

The second procedure cannot by itself establish new physics. An anomaly may
result from detector effects, calibration errors, statistical fluctuations,
or deficiencies in the background model. Nevertheless, it provides a powerful
mechanism for discovering phenomena not represented in a predefined catalogue
of theories.

\section{AI for Inverse Problems}

Many important problems in physics have the structure
\begin{equation}
\mathcal{F}[u]=y,
\end{equation}
where $u$ is an unknown physical field or parameter and $y$ represents
observations.

The forward problem is
\begin{equation}
u\longrightarrow y,
\end{equation}
while the inverse problem is
\begin{equation}
y\longrightarrow u.
\end{equation}

Inverse problems are frequently ill-posed in the sense of Hadamard. Small
errors in $y$ can produce large errors in the reconstructed $u$.

Machine learning can provide an approximate inverse map
\begin{equation}
u\simeq \mathcal{N}_\theta(y).
\end{equation}

Examples include:

\begin{itemize}
\item reconstruction of particle trajectories;
\item tomography;
\item reconstruction of cosmological parameters;
\item imaging of astrophysical objects;
\item estimation of material properties;
\item reconstruction of electromagnetic fields;
\item parameter inference in quantum systems.
\end{itemize}

However, an AI reconstruction should not be interpreted as a unique physical
solution merely because it produces a visually convincing result. The
non-uniqueness of the inverse problem remains a property of the underlying
physics.

\section{Physics-Informed Machine Learning}

One of the most important developments is the attempt to incorporate physical
laws directly into the learning procedure.

Suppose a physical field satisfies
\begin{equation}
\mathcal{L}[u]=0.
\end{equation}

A neural network proposes
\begin{equation}
u_\theta(x,t).
\end{equation}

Instead of minimizing only a data error,
\begin{equation}
L_{\mathrm{data}} =
\sum_i
\left|u_\theta(x_i,t_i)-u_i\right|^2,
\end{equation}
one may introduce a physics loss
\begin{equation}
L_{\mathrm{phys}}=
\int
\left|
\mathcal{L}[u_\theta]
\right|^2,dx,dt.
\end{equation}

The total loss is then
\begin{equation}
L=
\lambda_{\mathrm{data}}L_{\mathrm{data}}
+
\lambda_{\mathrm{phys}}L_{\mathrm{phys}}.
\end{equation}

This idea is commonly associated with physics-informed neural networks
(PINNs), although the broader concept includes many architectures and methods
\cite{Karn2021}.

The conceptual importance of this construction is greater than the particular
neural-network implementation. It represents an attempt to combine two sources
of knowledge:
\begin{equation}
\boxed{
\text{empirical information}
+
\text{theoretical information}
}.
\end{equation}

Physics is therefore not treated merely as a source of training examples.

\section{Exact Constraints and Symmetry}

A major weakness of unconstrained machine learning is that a network may
approximate a function extremely well while violating an exact physical
principle.

For example, a model may violate:

\begin{itemize}
\item conservation of energy;
\item momentum conservation;
\item gauge invariance;
\item Lorentz invariance;
\item permutation symmetry;
\item rotational symmetry;
\item unitarity;
\item positivity;
\item locality;
\item dimensional consistency.
\end{itemize}

Whenever possible, these properties should be built into the representation.

Suppose a physical quantity satisfies
\begin{equation}
f(Rx)=f(x)
\end{equation}
under a symmetry group $G$.

Rather than learning this symmetry from data, one can construct an architecture
that satisfies it exactly.

More generally, if
\begin{equation}
x\mapsto \rho_X(g)x,
\qquad
y\mapsto \rho_Y(g)y,
\end{equation}
then an equivariant map satisfies
\begin{equation}
f(\rho_X(g)x)=
\rho_Y(g)f(x).
\end{equation}

Equivariant neural networks provide a natural bridge between group theory and
machine learning.

This observation is particularly important for physics. Symmetry is not merely
a statistical regularity. It is often part of the definition of the physical
theory.

\section{Neural Operators and the Learning of Equations}

A conventional neural network usually learns a finite-dimensional map
\begin{equation}
\mathbb{R}^n\rightarrow\mathbb{R}^m.
\end{equation}

Many physical problems, however, are naturally formulated as mappings between
functions:
\begin{equation}
u(x)\longrightarrow v(x).
\end{equation}

A differential equation defines precisely such a mapping. For example,
\begin{equation}
\mathcal{L}*\mu u=f
\end{equation}
defines an operator
\begin{equation}
\mathcal{G}*\mu:f\mapsto u.
\end{equation}

Neural operators attempt to learn such mappings directly:
\begin{equation}
\mathcal{G}_\theta\simeq\mathcal{G}.
\end{equation}

This can be particularly valuable when the same physical problem must be
solved repeatedly for many parameter values. Once the operator has been
learned, new solutions can be obtained much more rapidly than by repeatedly
solving the original equations numerically
\cite{Azizzadenesheli2024}.

This suggests an important change in computational physics:

\begin{equation}
\text{solve equation repeatedly}
\quad\longrightarrow\quad
\text{learn the solution operator}.
\end{equation}

The latter does not eliminate the original physics. Rather, it constructs a
computational surrogate for it.

\section{AI and Numerical Physics}

There is a common misconception that machine learning is intended to replace
numerical physics.

In reality, the most effective strategy is often hybrid.

Suppose a computationally expensive simulation produces
\begin{equation}
y=F(x).
\end{equation}

A neural network can approximate
\begin{equation}
F(x)\simeq F_\theta(x).
\end{equation}

The expensive solver is used to generate a training set, while the learned
model is subsequently used for rapid exploration.

This is especially useful when one needs millions of evaluations, for example
in:

\begin{itemize}
\item Bayesian parameter estimation;
\item uncertainty propagation;
\item optimization;
\item Monte Carlo sampling;
\item experimental design;
\item real-time control.
\end{itemize}

Thus the appropriate comparison is not
\begin{equation}
\text{AI versus numerical physics},
\end{equation}
but
\begin{equation}
\boxed{
\text{AI + numerical physics}
}.
\end{equation}

\section{AI in Fluid Dynamics and Complex Systems}

Fluid dynamics provides an especially natural domain for AI because realistic
flows involve nonlinear partial differential equations and enormous ranges of
spatial and temporal scales.

The Navier--Stokes equations,
\begin{equation}
\partial_t\bm{u}
+
(\bm{u}\cdot\nabla)\bm{u}=
-\frac{1}{\rho}\nabla p
+
\nu\nabla^2\bm{u}
+
\bm{f},
\end{equation}
are computationally expensive at high Reynolds number.

Machine learning can be used to approximate unresolved scales, construct
reduced-order models, estimate flow fields from sparse measurements, and
control flows \cite{Brunton2020,Vinuesa2023}.

The important methodological possibility is the construction of hybrid models,
schematically
\begin{equation}
\partial_t u=
F_{\mathrm{known}}[u]
+
F_{\mathrm{ML}}[u].
\end{equation}

Here the well-established part of the physics remains explicit, while AI
represents an unknown or computationally expensive component.

This approach may be more scientifically robust than asking a neural network
to learn the entire dynamics from data.

\section{AI in Cosmology and Astrophysics}

Cosmology combines enormous datasets with expensive forward models, making it
particularly suitable for machine learning.

Possible applications include

\begin{itemize}
\item classification of galaxies;
\item gravitational-lensing reconstruction;
\item parameter estimation;
\item emulation of cosmological simulations;
\item detection of transient astronomical events;
\item reconstruction of large-scale structure;
\item inference of dark-matter distributions.
\end{itemize}

An important example is the use of learned surrogate models for cosmological
simulations. A full simulation may require substantial computational resources,
while a trained emulator can provide approximate predictions at a fraction of
the cost.

The danger is extrapolation. If the training simulations cover only a limited
region of parameter space, the learned model may produce plausible but
incorrect predictions elsewhere.

Consequently, the question
\begin{equation}
\text{`How accurate is the neural network?'}
\end{equation}
must be replaced by the more physical question
\begin{equation}
\text{`In what domain is the learned approximation valid?'}
\end{equation}

\section{AI in Quantum Physics}

Machine learning has found applications in quantum physics at several levels.

For a quantum state
\begin{equation}
|\psi\rangle=
\sum_x \psi(x)|x\rangle,
\end{equation}
a neural-network representation may approximate the amplitudes
\begin{equation}
\psi(x)\simeq\psi_\theta(x).
\end{equation}

Such representations have been used for variational calculations and many-body
quantum states \cite{Carleo2017,Carleo2019}.

Machine learning can also assist in:

\begin{itemize}
\item quantum state tomography;
\item quantum control;
\item identification of quantum phases;
\item many-body simulation;
\item Hamiltonian learning;
\item lattice field theory;
\item Monte Carlo sampling.
\end{itemize}

A particularly interesting possibility is learning a Hamiltonian from
observations. If measurements $D$ are generated by
\begin{equation}
H(\lambda)=\sum_i\lambda_i O_i,
\end{equation}
one may attempt to infer the parameters $\lambda_i$.

This is an inverse problem, but it also approaches the boundary between
parameter estimation and theory identification.

\section{AI and Fundamental Physics}

The most ambitious application of AI is the search for new fundamental physics.

The traditional scientific cycle is approximately
\begin{equation}
\text{observation}
\rightarrow
\text{concept}
\rightarrow
\text{mathematical theory}
\rightarrow
\text{prediction}
\rightarrow
\text{experiment}.
\end{equation}

AI potentially modifies this cycle by introducing a computational exploration
stage:
\begin{equation}
\text{observation}
\rightarrow
\boxed{\text{AI exploration}}
\rightarrow
\text{candidate structure}
\rightarrow
\text{human formulation}
\rightarrow
\text{prediction}.
\end{equation}

Symbolic regression is particularly relevant. Given observations
$(x_i,y_i)$, one searches over mathematical expressions
\begin{equation}
y=f(x)
\end{equation}
rather than fixing the functional form in advance.

The objective may be
\begin{equation}
\min_f
\left[
\chi^2(f)+\lambda C(f)
\right],
\end{equation}
where $C(f)$ measures complexity.

This expresses an old scientific principle in algorithmic form:

\begin{quote}
Prefer a simple law that explains the observations.
\end{quote}

However, finding a compact formula is not equivalent to discovering a physical
law. A law must possess explanatory power, survive independent tests, and fit
within a broader theoretical structure.

\section{AI and the Discovery of Physical Laws}

Suppose a set of observations suggests
\begin{equation}
y=f(x_1,x_2,\ldots,x_n).
\end{equation}

A machine-learning system may discover a highly accurate approximation
\begin{equation}
f_{\mathrm{ML}}(x_1,\ldots,x_n).
\end{equation}

But physics asks a deeper question:

\begin{equation}
\text{Why does this relation hold?}
\end{equation}

For example, the discovery of
\begin{equation}
F\propto\frac{m_1m_2}{r^2}
\end{equation}
is not equivalent to fitting an inverse-square function to observations. The
physical significance comes from its relation to a wider theory involving
symmetry, dynamics, conservation laws, and limiting principles.

This distinction can be summarized as

\begin{equation}
\boxed{
\text{prediction}\neq\text{explanation}
}.
\end{equation}

AI can potentially discover regularities. Whether those regularities constitute
physical understanding remains a separate question.

\section{The Problem of Interpretability}

Most deep neural networks contain a very large number of parameters. Their
internal representation is therefore difficult to interpret in conventional
physical terms.

The problem is not merely philosophical.

A black-box model may produce
\begin{equation}
y_{\mathrm{ML}}\approx y_{\mathrm{experiment}}
\end{equation}
with extraordinary accuracy while relying on correlations that have no physical
meaning.

For physics, interpretability can therefore be regarded as a scientific
requirement rather than merely a desirable software feature.

Several strategies are possible:

\begin{enumerate}
\item symbolic regression;
\item sparse model discovery;
\item dimensional analysis;
\item symmetry constraints;
\item interpretable latent variables;
\item causal modelling;
\item uncertainty quantification;
\item comparison with known limiting cases.
\end{enumerate}

The ideal scientific model should not merely predict the data. It should expose
the structure responsible for the prediction.

\section{Dimensional Analysis as a Constraint on AI}

Dimensional analysis provides a particularly simple and powerful example of
physical prior knowledge.

Suppose
\begin{equation}
Q=Q(x_1,\ldots,x_n).
\end{equation}

The dimensions must satisfy
\begin{equation}
[Q]=
\prod_i[x_i]^{a_i}.
\end{equation}

An unconstrained neural network generally does not know this.

A physically informed model can instead operate on dimensionless variables.
Using Buckingham's $\Pi$ theorem, one may write
\begin{equation}
\Pi_1=f(\Pi_2,\ldots,\Pi_k).
\end{equation}

This can drastically reduce the effective dimensionality of the learning
problem.

It also illustrates a broader principle:

\begin{equation}
\boxed{
\text{the more physical structure is known, the less AI has to learn}.
}
\end{equation}

This is one of the central methodological arguments for physics-informed AI.

\section{Symmetry as a Form of Data Compression}

Symmetry has another interpretation in machine learning.

If a system is invariant under a group $G$, then many apparently different
configurations represent the same physical situation.

Learning them independently is wasteful.

For example, rotational invariance means that
\begin{equation}
f(Rx)=f(x).
\end{equation}

Instead of training on all rotated configurations, the architecture can identify
them as equivalent.

Thus symmetry provides a form of data compression.

This is particularly significant in physics because the relevant symmetry groups
are often known a priori.

For relativistic systems, Lorentz symmetry requires
\begin{equation}
f(\Lambda x)=\rho(\Lambda)f(x),
\end{equation}
with $\Lambda\in SO(1,3)$ or the appropriate covering group.

Embedding such structures into machine-learning models connects AI directly with
the mathematical architecture of physical theories.

\section{AI and the Scientific Method}

It is useful to compare AI-assisted physics with the traditional scientific
method.

A simplified traditional cycle is
\begin{equation}
\boxed{
\text{Question}
\rightarrow
\text{Hypothesis}
\rightarrow
\text{Prediction}
\rightarrow
\text{Experiment}
\rightarrow
\text{Revision}.
}
\end{equation}

An AI-assisted cycle can contain additional loops:
\begin{equation}
\boxed{
\text{Question}
\rightarrow
\text{AI exploration}
\rightarrow
\text{candidate hypotheses}
\rightarrow
\text{simulation}
\rightarrow
\text{experiment}
\rightarrow
\text{AI updating}.
}
\end{equation}

The crucial point is that AI need not replace the scientific method. It may
increase the number and sophistication of hypotheses that can be explored.

This resembles the historical role of mathematics and computers. Neither
mathematics nor computation eliminated physical reasoning; they amplified it.

\section{AI as a Microscope for Theory Space}

A useful conceptual interpretation is that AI may function as a
\emph{microscope for theory space}.

Consider a space $\mathcal{T}$ of possible models:
\begin{equation}
T\in\mathcal{T}.
\end{equation}

A human physicist can explore only a very small part of this space. AI can
perform large-scale searches over parameterizations, functional forms,
architectures, or candidate effective theories.

The process may be represented as
\begin{equation}
\mathcal{T}
\overset{\mathrm{AI}}{\longrightarrow}
\mathcal{T}*{\mathrm{candidate}}
\overset{\mathrm{physics}}{\longrightarrow}
\mathcal{T}*{\mathrm{viable}}.
\end{equation}

The second arrow is essential.

AI can search the space, but physical reasoning must determine whether the
candidate structures are meaningful.

\section{AI and Effective Theories}

The framework of effective field theory provides an especially interesting
setting.

Suppose the low-energy Lagrangian is
\begin{equation}
\mathcal{L}_{\mathrm{eff}}=
\mathcal{L}_0
+
\sum_i
\frac{c_i}{\Lambda^{d_i-4}}
\mathcal{O}_i.
\end{equation}

AI can be used to identify combinations of operators that best describe data,
estimate Wilson coefficients, or search for deviations from a reference theory.

The important physical constraints remain:

\begin{itemize}
\item symmetry;
\item locality;
\item dimensional analysis;
\item unitarity;
\item gauge invariance;
\item consistency of the low-energy expansion.
\end{itemize}

Thus AI could become a search engine through the space of effective theories,
but it should not be allowed to treat mathematically possible expressions as
physically equivalent.

\section{Uncertainty Quantification}

One of the most serious weaknesses of ordinary machine learning is that
prediction accuracy does not automatically provide reliable uncertainty.

Physics requires statements such as
\begin{equation}
x=x_0\pm\Delta x,
\end{equation}
not merely
\begin{equation}
x\simeq x_0.
\end{equation}

For a learned prediction
\begin{equation}
y=f_\theta(x),
\end{equation}
one should ideally determine
\begin{equation}
p(y|x,D),
\end{equation}
where $D$ is the training dataset.

Uncertainty has several sources:

\begin{enumerate}
\item measurement noise;
\item finite training data;
\item model uncertainty;
\item numerical approximation;
\item distribution shift;
\item extrapolation beyond the training domain.
\end{enumerate}

The last category is particularly dangerous.

A model may be extremely accurate inside
\begin{equation}
x\in\mathcal{D}_{\mathrm{train}}
\end{equation}
and completely unreliable outside it.

This is especially problematic in physics because the most interesting discoveries
often occur precisely where existing models have not been tested.

\section{The Problem of Extrapolation}

Interpolation and extrapolation must be sharply distinguished.

If
\begin{equation}
x_{\min}<x<x_{\max}
\end{equation}
belongs to the training domain, the model may interpolate successfully.

For
\begin{equation}
x<x_{\min}
\quad\text{or}\quad
x>x_{\max},
\end{equation}
it extrapolates.

Physics frequently demands extrapolation:

\begin{itemize}
\item high-energy limits;
\item low-temperature limits;
\item early-universe conditions;
\item long-time evolution;
\item extreme densities;
\item strong gravitational fields.
\end{itemize}

A purely data-driven model has no general reason to extrapolate according to
the correct physical law.

Physical priors can therefore be more important for extrapolation than for
interpolation.

\section{AI and Causality}

Correlation is not causation.

Suppose
\begin{equation}
Y=f(X)
\end{equation}
is observed experimentally.

A machine-learning model may learn this relation without determining whether
$X$ causes $Y$, whether $Y$ causes $X$, or whether both are consequences of a
third variable $Z$.

Physics is fundamentally concerned with causal structure because dynamical
equations specify how physical states evolve.

A future scientific AI should therefore incorporate causal structure rather than
merely statistical dependence.

This is especially important in complex systems where correlations may be
strong but physically irrelevant.

\section{Generative AI and Large Language Models}

Large language models introduce a qualitatively different capability.

A physicist can ask an LLM to:

\begin{itemize}
\item explain a mathematical derivation;
\item generate computer code;
\item manipulate symbolic expressions;
\item search and organize literature;
\item compare theoretical approaches;
\item propose possible calculations;
\item translate between mathematical representations;
\item assist in writing scientific documents.
\end{itemize}

These capabilities can substantially reduce the cost of routine intellectual
work.

However, LLMs have an important limitation: they generate plausible symbolic
sequences rather than guaranteeing mathematical truth.

Consequently, an LLM-generated derivation must be treated like a conjecture
until verified.

This distinction is fundamental:
\begin{equation}
\boxed{
\text{linguistic plausibility}\neq\text{mathematical validity}.
}
\end{equation}

For physics, automatic verification through symbolic algebra, numerical tests,
dimensional analysis, limiting cases, and independent derivations should
therefore accompany AI-generated mathematics.

\section{AI as a Collaborator in Theoretical Physics}

A productive model of AI-assisted theory research is not
\begin{equation}
\text{human}\rightarrow\text{AI}\rightarrow\text{answer},
\end{equation}
but
\begin{equation}
\text{human}
\leftrightarrow
\text{AI}
\leftrightarrow
\text{computer algebra}
\leftrightarrow
\text{numerical simulation}
\leftrightarrow
\text{literature}.
\end{equation}

The human physicist remains responsible for:

\begin{itemize}
\item defining the physical problem;
\item selecting meaningful variables;
\item identifying relevant principles;
\item interpreting results;
\item checking consistency;
\item deciding what constitutes an explanation.
\end{itemize}

AI can contribute:

\begin{itemize}
\item rapid exploration;
\item pattern recognition;
\item generation of candidate models;
\item code generation;
\item optimization;
\item large-scale literature synthesis;
\item automated numerical experimentation.
\end{itemize}

This division of labor is likely to be more productive than attempting to
construct a completely autonomous scientist.

\section{Verification: The Essential Scientific Filter}

Any AI-generated physical result should pass several independent tests.

\subsection{Dimensional consistency}

Every equation must possess consistent physical dimensions.

\subsection{Symmetry}

The result must respect the relevant symmetry principles.

\subsection{Limiting cases}

For example,
\begin{equation}
\hbar\rightarrow0,
\qquad
c\rightarrow\infty,
\qquad
G\rightarrow0,
\end{equation}
should reproduce appropriate classical, nonrelativistic, or nongravitational
limits where applicable.

\subsection{Conservation laws}

If the theory predicts a conserved quantity,
\begin{equation}
\frac{dQ}{dt}=0,
\end{equation}
the learned model must reproduce this property.

\subsection{Independent reproduction}

The result should ideally be confirmed by a method independent of the original
AI system.

This last criterion is especially important because an AI system may reproduce
its own systematic errors.

\section{AI Does Not Eliminate the Need for Theory}

There is a tempting but misleading argument:

\begin{quote}
If a sufficiently powerful neural network can predict every observable, then a
physical theory is unnecessary.
\end{quote}

This conclusion does not follow.

A predictive algorithm and a physical theory have different functions.

A theory provides:

\begin{enumerate}
\item a compact description;
\item a conceptual framework;
\item a mechanism;
\item relations among apparently different phenomena;
\item extrapolation principles;
\item consistency conditions;
\item new predictions.
\end{enumerate}

An AI model may perform the first task extremely well while providing little
insight into the others.

The distinction is analogous to the difference between a table of experimental
results and a theory that explains them.

\section{The Role of Mathematical Physics}

Mathematical physics may become more important, rather than less important, in
the age of AI.

Machine learning produces candidate structures. Mathematics provides the tools
for determining whether those structures are:

\begin{itemize}
\item well-defined;
\item stable;
\item differentiable;
\item invariant;
\item causal;
\item consistent;
\item integrable;
\item unitary;
\item renormalizable, where appropriate.
\end{itemize}

For example, if AI proposes a new dynamical equation
\begin{equation}
\partial_t u=F[u],
\end{equation}
mathematical analysis can determine existence, uniqueness, stability, and
well-posedness.

Thus AI may increase the number of conjectures while mathematics remains
essential for establishing their validity.

\section{AI and the Future of Experimental Design}

AI can be particularly powerful when experiments are expensive.

Suppose an experiment depends on parameters
\begin{equation}
\bm{\lambda}=
(\lambda_1,\ldots,\lambda_n).
\end{equation}

Rather than selecting $\bm{\lambda}$ according to a fixed grid, AI can optimize
the experiment according to an information criterion.

For Bayesian inference one may maximize the expected information gain
\begin{equation}
\mathcal{I}=
D_{\mathrm{KL}}
\left[
p(\theta|D,\bm{\lambda})
\Vert
p(\theta)
\right].
\end{equation}

The experimental system can therefore become adaptive:
\begin{equation}
\text{measurement}
\rightarrow
\text{analysis}
\rightarrow
\text{updated model}
\rightarrow
\text{next measurement}.
\end{equation}

This creates a feedback loop between theory, computation, and experiment.

\section{Towards Autonomous Scientific Laboratories}

The logical extension is an autonomous scientific laboratory.

A simplified architecture would be

\begin{equation}
\boxed{
\begin{array}{c}
\text{AI hypothesis generation}\
\downarrow\
\text{simulation}\
\downarrow\
\text{experimental design}\
\downarrow\
\text{automated experiment}\
\downarrow\
\text{data analysis}\
\downarrow\
\text{hypothesis updating}
\end{array}}
\end{equation}

Such systems could explore regions of parameter space much more rapidly than
human researchers.

Nevertheless, complete autonomy should not be confused with scientific
understanding. An automated system may optimize a criterion without knowing why
the criterion is scientifically interesting.

The definition of the scientific objective remains a profoundly human task.

\section{A Taxonomy of AI Roles in Physics}

The discussion can be summarized by the following hierarchy.

\begin{center}
\begin{tabular}{lll}
\toprule
Level & AI role & Scientific function\\
\midrule
I & Calculator & Numerical acceleration\\
II & Classifier & Data interpretation\\
III & Emulator & Fast simulation\\
IV & Inference engine & Parameter reconstruction\\
V & Optimizer & Experimental design\\
VI & Pattern finder & Anomaly detection\\
VII & Model finder & Equation discovery\\
VIII & Hypothesis generator & Theory exploration\\
IX & Scientific agent & Closed-loop discovery\\
\bottomrule
\end{tabular}
\end{center}

The first five levels are already well established. Levels VI--VIII are active
areas of research. Level IX remains largely prospective.

\section{A Possible Division of Labor Between Physics and AI}

The most promising methodology can be summarized schematically as

\begin{equation}
\boxed{
\begin{array}{rcl}
\text{Physics} &\longrightarrow&
\text{principles, symmetries, constraints, questions}\\
&&\downarrow\\
\text{AI} &\longrightarrow&
\text{search, approximation, inference, optimization}\\
&&\downarrow\\
\text{Mathematics} &\longrightarrow&
\text{consistency, proof, structure}\\
&&\downarrow\\
\text{Experiment} &\longrightarrow&
\text{empirical validation}.
\end{array}}
\end{equation}

The essential idea is that none of these components is sufficient by itself.

Data without theory can produce correlations.

Theory without data can produce elegant but physically false models.

AI without mathematical verification can produce convincing errors.

Mathematics without empirical input cannot determine which possible structures
describe nature.

The future of physics is therefore likely to be increasingly hybrid.

\section{Limitations and Risks}

Several limitations deserve particular emphasis.

\subsection{Training-data bias}

AI learns from its data. If the training set is incomplete or biased, the
resulting model inherits those deficiencies.

\subsection{Distribution shift}

A model trained under one experimental condition may fail under another.

\subsection{Hallucination}

Generative AI can produce references, equations, or arguments that appear
credible but are false.

\subsection{Opacity}

A high-dimensional model may be difficult to interpret physically.

\subsection{Overfitting}

A model can fit statistical fluctuations rather than genuine physical structure.

\subsection{False discovery}

Given sufficiently large datasets and sufficiently flexible models, apparently
significant correlations can arise by chance.

\subsection{Reproducibility}

Scientific AI requires reproducible datasets, model architectures, training
procedures, random seeds where relevant, and validation protocols.

These issues mean that AI must be embedded within the existing culture of
scientific verification rather than treated as an oracle.

\section{The Deeper Philosophical Question}

The use of AI in physics raises a question that is deeper than technology:

\begin{quote}
What does it mean to understand a physical phenomenon?
\end{quote}

Suppose an AI system predicts every experimentally accessible quantity with
extraordinary accuracy.

Does it understand the theory?

There are at least three possible meanings of understanding:

\begin{enumerate}
\item \textbf{Predictive understanding:} the ability to predict observations.

\item \textbf{Structural understanding:} identification of mathematical
relations and symmetries.

\item \textbf{Conceptual understanding:} identification of the physical
principles and mechanisms responsible for the phenomena.
\end{enumerate}

AI is already very powerful in the first category and increasingly capable in
the second.

The third remains much less clear.

This distinction may become central to the philosophy of AI-assisted science.

\section{Conclusion}

Artificial intelligence is becoming an integral component of physics. Its role
already extends from experimental data analysis to numerical simulation,
inverse problems, experimental design, and theoretical research.

The most important development is not the replacement of physical equations by
neural networks. It is the emergence of hybrid methods in which data-driven and
theory-driven approaches are combined.

Physics supplies structure:
\begin{equation}
\text{symmetry}+\text{conservation}+\text{dynamics}+\text{causality}.
\end{equation}

AI supplies computational flexibility:
\begin{equation}
\text{search}+\text{approximation}+\text{pattern recognition}
+\text{optimization}.
\end{equation}

The resulting methodology is potentially more powerful than either component
alone:
\begin{equation}
\boxed{
\text{AI}+\text{physics}+\text{mathematics}+\text{experiment}.
}
\end{equation}

The greatest opportunity may lie not in asking AI to discover physics without
human intervention, but in using AI to enlarge the domain of problems that
physicists can investigate.

The historical progression
\begin{equation}
\text{experiment}
\rightarrow
\text{theory}
\rightarrow
\text{computation}
\rightarrow
\text{data-intensive science}
\end{equation}
may therefore be followed by
\begin{equation}
\boxed{
\text{human--AI scientific collaboration}.
}
\end{equation}

Whether this will constitute a genuinely new paradigm of science depends not
on the size of neural networks but on whether AI systems can participate in the
central activities of physics: identifying meaningful concepts, discovering
general principles, producing falsifiable predictions, and revealing structures
that humans did not previously recognize.

For the foreseeable future, the most scientifically productive vision is not
``AI instead of physics'', but

\begin{equation}
\boxed{
\text{AI as an instrument for doing physics that was previously
computationally or conceptually inaccessible.}
}
\end{equation}

\begin{thebibliography}{99}

\bibitem{Carleo2017}
G.~Carleo and M.~Troyer,
``Solving the quantum many-body problem with artificial neural networks,''
Science \textbf{355}, 602--606 (2017).

\bibitem{Carleo2019}
G.~Carleo, I.~Cirac, K.~Cranmer, L.~Daudet, M.~Schuld, N.~Tishby,
L.~Vogt-Maranto and L.~Zdeborov'a,
``Machine learning and the physical sciences,''
Rev.\ Mod.\ Phys.\ \textbf{91}, 045002 (2019).

\bibitem{Brunton2020}
S.~L.~Brunton, B.~R.~Noack and P.~Koumoutsakos,
``Machine learning for fluid mechanics,''
Annu.\ Rev.\ Fluid Mech.\ \textbf{52}, 477--508 (2020).

\bibitem{Karn2021}
G.~E.~Karniadakis, I.~G.~Kevrekidis, L.~Lu, P.~Perdikaris,
S.~Wang and L.~Yang,
``Physics-informed machine learning,''
Nature Rev.\ Phys.\ \textbf{3}, 422--440 (2021).

\bibitem{Karagiorgi2022}
G.~Karagiorgi, G.~Kasieczka, S.~Kravitz, B.~Nachman and D.~Shih,
``Machine learning in the search for new fundamental physics,''
Nature Rev.\ Phys.\ \textbf{4}, 399--412 (2022).

\bibitem{Douglas2022}
M.~R.~Douglas,
``Machine learning as a tool in theoretical science,''
Nature Rev.\ Phys.\ \textbf{4}, 145--146 (2022).

\bibitem{Thiyagalingam2022}
J.~Thiyagalingam, M.~Shankar, G.~Fox and T.~Hey,
``Scientific machine learning benchmarks,''
Nature Rev.\ Phys.\ \textbf{4}, 413--420 (2022).

\bibitem{Cranmer2023}
K.~Cranmer, G.~Kanwar, S.~Racani\`ere, D.~J.~Rezende and
P.~E.~Shanahan,
``Advances in machine-learning-based sampling motivated by lattice quantum
chromodynamics,''
Nature Rev.\ Phys.\ \textbf{5}, 526--535 (2023).

\bibitem{Vinuesa2023}
R.~Vinuesa, S.~L.~Brunton and B.~McKeon,
``The transformative potential of machine learning for experiments in fluid
mechanics,''
Nature Rev.\ Phys.\ \textbf{5}, 536--545 (2023).

\bibitem{Azizzadenesheli2024}
K.~Azizzadenesheli, Z.~Kovachki, Z.~Li, M.~Liu-Schiaffini,
J.~Kossaifi and A.~Anandkumar,
``Neural operators for accelerating scientific simulations and design,''
Nature Rev.\ Phys.\ \textbf{6}, 320--328 (2024).

\bibitem{He2024}
Y.-H.~He,
``AI-driven research in pure mathematics and theoretical physics,''
Nature Rev.\ Phys.\ \textbf{6}, 546--553 (2024).

\bibitem{Bracco2025}
A.~Bracco, J.~Brajard, H.~A.~Dijkstra, P.~Hassanzadeh,
C.~Lessig and C.~Monteleoni,
``Machine learning for the physics of climate,''
Nature Rev.\ Phys.\ \textbf{7}, 6--20 (2025).

\end{thebibliography}

\end{document}
